\section{Methodology}\label{sec:method}

\subsection{Host and Software Stack}
All numbers in this paper come from a single certification run on a
host fingerprinted as \emph{xbabe1}: 128 vCPU x86\_64, Linux 6.8 with
ext4, Docker 29.2.1, Rust 1.95.0, and SQLite 3.x bundled by the
\texttt{rusqlite} 0.37 crate~\cite{rusqlite}. The image digest, host
fingerprint, and per-output checksums are recorded in
\texttt{manifest.json}; the appendix excerpts the relevant fields.

\subsection{Workload Matrix}
The certification matrix is the Cartesian product of 8 workloads, 2
durability modes (strict and relaxed), 8 thread levels
($\{1, 2, 4, 8, 16, 32, 64, 128\}$), 5 measured repetitions, and 2
engines (RedlineDB and SQLite). One warmup repetition per cell is run
and discarded. The grand total is $1{,}728$ child processes:
$1{,}280$ measured and $448$ warmup. Every measured run is one fresh
child process with its own database directory, so no harness state
crosses between runs. The eight workloads cover the OLTP design
space: \texttt{point-read-pk}, \texttt{point-read-secondary},
\texttt{secondary-index-range}, \texttt{mixed95-read5-write},
\texttt{mixed50-read50-write}, \texttt{writers-disjoint},
\texttt{hot-row-update}, and \texttt{insert-only}. The shapes are
inspired by the standard OLTP benchmark
literature---TPC-C~\cite{tpcc}, YCSB~\cite{cooper2010ycsb},
OLTP-Bench~\cite{difallah2013oltpbench}, and
sysbench~\cite{sysbench}---without copying any one harness
end-to-end, since the certification protocol's bookkeeping (real
warmup discard, recursive on-disk footprint, BUSY/LOCKED/timeout
split metrics, fsynced ack log oracle) is not standard in those
suites.

\subsection{Per-Run Telemetry}
Each child process emits a JSONL row containing throughput
(operations/second), latency percentiles
(p50/p95/p99/p999/max in microseconds), BUSY/LOCKED/timeout split
error counts, peak RSS, fdatasync and pwrite syscall counts (passed
through from the kernel's WAL counters and cross-checked with
\texttt{strace} on Linux), recursive on-disk \texttt{data\_bytes},
an integrity checksum, and the host fingerprint. Across the 1{,}280
measured runs we observed 0 hard failures.

\subsection{Recovery Matrix}
The recovery matrix is 36 cases: 3 crash scenarios (clean shutdown,
mid-WAL crash, post-WAL pre-commit crash) crossed with 2 durability
modes and 6 repetitions each. Each case writes a known transaction
prefix, kills the writer at a deterministic point, restarts, and
verifies the post-recovery state matches the prefix the oracle has
acked. Pass count: 36/36.

\subsection{Failpoint Matrix}
The failpoint matrix is 24 cases: a deterministic fail-point hit
schedule across 7 sites
(\texttt{wal::write\_encoded}, \texttt{wal::flush},
\texttt{wal::flush\_until}, \texttt{wal::flush\_all},
\texttt{commit::publish}, \texttt{heap::mut}, \texttt{index::mut},
\texttt{catalog::rename}) under strict durability. Each case panics
the child at a controlled hit count, restarts via the same recovery
path, and the oracle confirms zero lost-acked-commits. Pass count:
24/24.

\subsection{SQL Compatibility Matrix}
The compatibility matrix is 40 SQLLogicTest-style
cases~\cite{sqllogictest} drawn from the SQLite-dialect surface
RedlineDB targets. The cases exercise CREATE/SELECT/INSERT/UPDATE/
DELETE/INDEX/TRANSACTION semantics, NULL handling, and the
five-class type affinity rules where they remain user-visible.
Pass count: 40/40.

\subsection{Threats to Validity}
The single-host caveat is real: every number in this paper comes
from one machine. We mitigate it by reporting the host fingerprint
verbatim and by tagging the certification commit so any reader with
similar hardware can rerun. The SQLite version we compare against is
the one bundled by \texttt{rusqlite} 0.37, not a hand-built
amalgamation, so we are measuring the shipping
ecosystem~\cite{rusqlite} rather than a custom build. We do not
benchmark against libSQL~\cite{libsql} or DuckDB~\cite{raasveldt2019duckdb}
in this submission; libSQL's WAL contract is the same as upstream
SQLite by construction, and DuckDB's design center is column-store
analytics rather than OLTP, so neither is the right comparison
target for the specific concurrency-cliff hypothesis we test.
